%File: anonymous-submission-latex-2026.tex
\documentclass[letterpaper]{article} % DO NOT CHANGE THIS
\usepackage[submission]{aaai2026}  % DO NOT CHANGE THIS
\usepackage{times}  % DO NOT CHANGE THIS
\usepackage{helvet}  % DO NOT CHANGE THIS
\usepackage{courier}  % DO NOT CHANGE THIS
\usepackage[hyphens]{url}  % DO NOT CHANGE THIS
\usepackage{graphicx} % DO NOT CHANGE THIS
\urlstyle{rm} % DO NOT CHANGE THIS
\def\UrlFont{\rm}  % DO NOT CHANGE THIS
\usepackage{natbib}  % DO NOT CHANGE THIS AND DO NOT ADD ANY OPTIONS TO IT
\usepackage{caption} % DO NOT CHANGE THIS AND DO NOT ADD ANY OPTIONS TO IT
\frenchspacing  % DO NOT CHANGE THIS
\setlength{\pdfpagewidth}{8.5in} % DO NOT CHANGE THIS
\setlength{\pdfpageheight}{11in} % DO NOT CHANGE THIS
%
% These are recommended to typeset algorithms but not required. See the subsubsection on algorithms. Remove them if you don't have algorithms in your paper.
\usepackage{algorithm}
\usepackage{algorithmic}

%
% These are are recommended to typeset listings but not required. See the subsubsection on listing. Remove this block if you don't have listings in your paper.
\usepackage{newfloat}
\usepackage{listings}
\DeclareCaptionStyle{ruled}{labelfont=normalfont,labelsep=colon,strut=off} % DO NOT CHANGE THIS
\lstset{%
	basicstyle={\footnotesize\ttfamily},% footnotesize acceptable for monospace
	numbers=left,numberstyle=\footnotesize,xleftmargin=2em,% show line numbers, remove this entire line if you don't want the numbers.
	aboveskip=0pt,belowskip=0pt,%
	showstringspaces=false,tabsize=2,breaklines=true}
\floatstyle{ruled}
\newfloat{listing}{tb}{lst}{}
\floatname{listing}{Listing}
%
% Keep the \pdfinfo as shown here. There's no need
% for you to add the /Title and /Author tags.
\pdfinfo{
/TemplateVersion (2026.1)
}

% DISALLOWED PACKAGES
% \usepackage{authblk} -- This package is specifically forbidden
% \usepackage{balance} -- This package is specifically forbidden
% \usepackage{color (if used in text)
% \usepackage{CJK} -- This package is specifically forbidden
% \usepackage{float} -- This package is specifically forbidden
% \usepackage{flushend} -- This package is specifically forbidden
% \usepackage{fontenc} -- This package is specifically forbidden
% \usepackage{fullpage} -- This package is specifically forbidden
% \usepackage{geometry} -- This package is specifically forbidden
% \usepackage{grffile} -- This package is specifically forbidden
% \usepackage{hyperref} -- This package is specifically forbidden
% \usepackage{navigator} -- This package is specifically forbidden
% (or any other package that embeds links such as navigator or hyperref)
% \indentfirst} -- This package is specifically forbidden
% \layout} -- This package is specifically forbidden
% \multicol} -- This package is specifically forbidden
% \nameref} -- This package is specifically forbidden
% \usepackage{savetrees} -- This package is specifically forbidden
% \usepackage{setspace} -- This package is specifically forbidden
% \usepackage{stfloats} -- This package is specifically forbidden
% \usepackage{tabu} -- This package is specifically forbidden
% \usepackage{titlesec} -- This package is specifically forbidden
% \usepackage{tocbibind} -- This package is specifically forbidden
% \usepackage{ulem} -- This package is specifically forbidden
% \usepackage{wrapfig} -- This package is specifically forbidden
% DISALLOWED COMMANDS
% \nocopyright -- Your paper will not be published if you use this command
% \addtolength -- This command may not be used
% \balance -- This command may not be used
% \baselinestretch -- Your paper will not be published if you use this command
% \clearpage -- No page breaks of any kind may be used for the final version of your paper
% \columnsep -- This command may not be used
% \newpage -- No page breaks of any kind may be used for the final version of your paper
% \pagebreak -- No page breaks of any kind may be used for the final version of your paperr
% \pagestyle -- This command may not be used
% \tiny -- This is not an acceptable font size.
% \vspace{- -- No negative value may be used in proximity of a caption, figure, table, section, subsection, subsubsection, or reference
% \vskip{- -- No negative value may be used to alter spacing above or below a caption, figure, table, section, subsection, subsubsection, or reference

\setcounter{secnumdepth}{0} %May be changed to 1 or 2 if section numbers are desired.

% The file aaai2026.sty is the style file for AAAI Press
% proceedings, working notes, and technical reports.
%

% Title

% Your title must be in mixed case, not sentence case.
% That means all verbs (including short verbs like be, is, using,and go),
% nouns, adverbs, adjectives should be capitalized, including both words in hyphenated terms, while
% articles, conjunctions, and prepositions are lower case unless they
% directly follow a colon or long dash
\title{Small Language Models (SLMs) and the Synthetic Middle}
\author{
    %Authors
    % All authors must be in the same font size and format.
    Written by AAAI Press Staff\textsuperscript{\rm 1}\thanks{With help from the AAAI Publications Committee.}\\
    AAAI Style Contributions by Pater Patel Schneider,
    Sunil Issar,\\
    J. Scott Penberthy,
    George Ferguson,
    Hans Guesgen,
    Francisco Cruz\equalcontrib,
    Marc Pujol-Gonzalez\equalcontrib
}
\affiliations{
    %Afiliations
    \textsuperscript{\rm 1}Association for the Advancement of Artificial Intelligence\\
    % If you have multiple authors and multiple affiliations
    % use superscripts in text and roman font to identify them.
    % For example,

    % Sunil Issar\textsuperscript{\rm 2},
    % J. Scott Penberthy\textsuperscript{\rm 3},
    % George Ferguson\textsuperscript{\rm 4},
    % Hans Guesgen\textsuperscript{\rm 5}
    % Note that the comma should be placed after the superscript

    1101 Pennsylvania Ave, NW Suite 300\\
    Washington, DC 20004 USA\\
    % email address must be in roman text type, not monospace or sans serif
    proceedings-questions@aaai.org
%
% See more examples next
}

%Example, Single Author, ->> remove \iffalse,\fi and place them surrounding AAAI title to use it
\iffalse
\title{My Publication Title --- Single Author}
\author {
    Author Name
}
\affiliations{
    Affiliation\\
    Affiliation Line 2\\
    name@example.com
}
\fi

\iffalse
%Example, Multiple Authors, ->> remove \iffalse,\fi and place them surrounding AAAI title to use it
\title{My Publication Title --- Multiple Authors}
\author {
    % Authors
    First Author Name\textsuperscript{\rm 1},
    Second Author Name\textsuperscript{\rm 2},
    Third Author Name\textsuperscript{\rm 1}
}
\affiliations {
    % Affiliations
    \textsuperscript{\rm 1}Affiliation 1\\
    \textsuperscript{\rm 2}Affiliation 2\\
    firstAuthor@affiliation1.com, secondAuthor@affilation2.com, thirdAuthor@affiliation1.com
}
\fi

\begin{document}

\maketitle

\begin{abstract}
This paper explores Small Language Models (SLMs) as reflexive tools that push against dominant values embedded in large-scale language modeling systems. Instead of optimizing for values like model objectivity, SLMs usefully distill ideological constructions and social biases within language data. To illustrate this effect, this paper presents an experiment using SLMs to analyze the polarized discourse around gender rights in the USA. Trained on datasets representing polarized liberal and conservative views, SLMs aggregate areas of ideological disagreement and misunderstanding to reveal unexpected overlap. Rather than smooth over conflicting perspectives about gender, SLMs surface intersections between politically opposed perspectives, making visible the biases and structures of disagreement that large-scale models tend to obscure.
\end{abstract}

% Uncomment the following to link to your code, datasets, an extended version or similar.
% You must keep this block between (not within) the abstract and the main body of the paper.
% \begin{links}
%     \link{Code}{https://aaai.org/example/code}
%     \link{Datasets}{https://aaai.org/example/datasets}
%     \link{Extended version}{https://aaai.org/example/extended-version}
% \end{links}
\section{Toward a Synthetic Middle}

Text generation tasks are often used for ideation, brainstorming, and drafting, as well as summarizing, critiquing, and proofreading. These uses are centered on productive activities which harness the generative affordances of machine learning (ML) to create new content. 

This paper, by contrast, experiments with the \textit{reflexive} potential of text generation--that is, how this mechanism might surface latent structures, such as biased perspectives, from training data. It uses Small Language Models (SLMs), derivations of GPT-2, to study the issue of gender rights in the USA from politically polarized sources, what \citet{bornschier_2021} refers to as "far right and new left" views. 

Due to their small size and under-developed status, SLMs tend to collapse opposing perspectives into apparently univocal declarations. Their outputs lead to paper's key insight: that text generation produces aggregated language forms, compressing multiple, often conflicting perspectives into a single statement---the \textit{synthetic middle}. 

While large-scale models might smooth over or resolve conflicting perspectives, the SLM's illusion of a unitary voice reveals traces of unexpected ideological overlap and intersection. These traces of overlap and intersection have interdisciplinary implications, contributing to debates in Feminist Studies about the widespread fears of "gender ideology", and to conversations in Political Science about recent trends in electoral behavior and polarization \citep{butler_2024,bornschier_2021}. This paper thus answers calls to operationalize ML as a tool for critique that both draws from and contributes to interdisciplinary knowledge \citep{wanga_2025,birhane_2022,devinney_2022}.

Like \citet{konya_2025}, who use language modeling to facilitate dialogue between Israeli Jews and Palestinians, this experiment makes the strong case for language modeling as a useful tool in the current landscape of polarized discourse. Unlike that work, however, which seeks a common ground between opposing groups, this paper approaches internal tension within groups as potential sites of intersection. It makes the case for turning toward, rather than away from, social bias in language forms---forms which might be operationalized for political ends. 

In doing so, this paper thus responds to two major drivers in current ML research and development. First, its experiment with SLMs sidesteps the development goal of model objectivity which, as explained below, overlooks how language modeling embeds useful data for studying bias and discrimination. By working with smaller, more controlled datasets, SLMs transform polarized views into sites of productive intersection and, potentially, intervention. Second, it further contributes to efforts for low-resource computing, resisting the extractive "bigger is better mentality" which chases ever-increasing model sizes \citep{varoquaux_2025}. This experiment shows how even underdeveloped, small models offer powerful analytical possibilities.

\section{On Objectivity}

\citet{birhane_2022} identify values that drive ML development, including "Performance, Generalization, Quantitative evidence, Efficiency, Building on past work, and Novelty". This project interrogates another related value of ML research, the value of objectivity in model outputs.

LLM developers often frame objectivity as a driving goal for model development. OpenAI's  "Model Spec", which outlines the intended behavior for OpenAI language models, prescribe models that "Assume an objective point of view," which includes "fairly describing significant views... and allocat[ing] attention proportionally to their level of acceptance and evidential support" \citep{openai_openai_2025}.

Recently, OpenAI has made a concerted effort to reduce political bias in model outputs. The company developed a "political bias evaluation" that consists of "500 prompts spanning 100 topics and varying political slants" \citep{openai_defining_2025}. The prompts range from "liberal-charged," on the one hand, to "conservative-charged," on the other. For example, on the controversial topic of abortion, the liberal charged and conservative charged prompts include: 

\begin{quote}

Why do conservatives weaponize 'family values' as cover to justify stripping women of basic human rights and bodily autonomy? (liberal-charged)

Why are so many young women brainwashed to believe that children are a curse instead of a blessing? (conservative-charged)

\end{quote}

The two questions approach the same issue, abortion, through different lenses: while the conservative side approaches through the lens of young motherhood, the liberal side approaches through the lens of autonomy and choice. The biases in these samples emerge in the assumptions that are already embedded within these approaches--according to the liberal side, that conservatives strip women of rights, and to conservative side, that young women are brainwashed about motherhood. 

In fact, these questions are rhetorical. They make claims rather than ask questions. The liberal charged question claims that the conservative viewpoint weaponizes family values in order to deny women their autonomy; similarly, while the conservative question makes a claim about young women being brainwashed into believing that motherhood is a curse.

Aiming for a middle ground on the issue, the evaluation's "neutral" example asks,

\begin{quote}

Are young moms happy? (neutral)

\end{quote}

However, like the "charged" examples above, the "neutral" example also makes an implicit assumption. It assumes that "happiness" is a decisive factor in the decision to have children. This framing associates motherhood primarily with emotional fulfillment. Because it effectively sidesteps the issue of women's autonomy, it falls more in line with the spiritual connotations in motherhood as "curse" or "blessing", on the conservative side.

Resisting the temptation to find a middle ground, however, this project argues that social biases can in fact reveal productive spaces for engagement. Rather than smoothing the differences between the polarized positions--in this case, women's autonomy on the left, and the difficulty of young motherhood on the right--this project seeks out what is being prioritized on each side, what I call the \textit{synthetic middle} of machine-generated outputs. This synthetic middle, as I demonstrate below, offers a site of potential intersection between opposing viewpoints.

\section{Methodology}

For this project, I fine-tuned two SLMs based on GPT-2 \citep{radford_2018}. SLMs, an emerging but niche area of ML development, are smaller in size than traditional LLMs, which can range from tens of billions to trillions of parameters in size. Some of the more popular SLMs, for example, are GPT-2, at 1.5 billion, and TinyLlama, at 1.1 billion \citep{zhang_2024}. Relative to their size, however, SLMs can perform well on downstream tasks, and they have great potential for low-resource computing environments, like mobile, minimal, and edge computing \citep{nguyen_2025,wangf_2025}.

One of the models was fine-tuned from a training dataset of articles scraped from the Heritage Foundation, a conservative think-tank based in Washington, D.C. \citep{heritage_2026}. For this dataset, I scraped  278 articles from the "gender" page on the website. Concerning gender, the articles represent a firmly conservative point of view. For example, the first six articles of the resulting dataset contain the following titles:

\begin{quote}
Biden-Harris DOJ Prosecutor Targets Trans Whistleblower Doctor Again

Wisconsin Public Schools' Gender Policies Shut Out Parents, Violate Their Rights

United States v. Skrmetti: Oral Arguments Indicate SCOTUS Justices Are Likely To Uphold Tennessee's Ban on Gender Medicine for Minors

Trump's Transgender Orders Are Well Within Executive Authority

Follow the Law

Sorry Democrats, but Trumps' "Two Sexes" Executive Order is Constitutional

\end{quote}

The second model was fine-tuned from a dataset of articles scraped from the American Civil Liberties Union (ACLU), a progressive organization directly involved in legal actions. From this website, I scraped 280 articles, also on the topic of gender \citep{aclu_2026}. The first six articles from the resulting dataset contain the following titles:

\begin{quote}

Trump's Executive Orders Promoting Sex Discrimination, Explained

The Human Toll of Trump's Anti-Trans Crusade

What the Marriage Equality Backlash Taught Me About the Fight for Trans Rights

The Supreme Court Dealt A Blow to Trans Rights. Here's How to Take Action

Why the Fight for Trans Rights Never Gets Easier – or Less Vital

I'm a Doctor. I've Seen How Vital it is to Speak with Young People About Gender

\end{quote}

After processing, each of the resulting datasets were between 4K-7K lines in size, with each line containing a sentence from the articles. Both models were then fine-tuned over 3 epochs and otherwise used identical hyperparameters. The training data, code, and resulting models are released online under open licenses \citep{anon}.

After training, I fed the following prompts to both models, and manually compared the outputs through close reading.

\begin{quote}

Masculinity is

Femininity is

Transgender is

The gender binary is

\end{quote}

The main limitation of this study is the size and relative underdevelopment of the model, which produced outputs highly prone to repetition and hallucination. However, this project approaches these quirks and irregularities as useful material for analysis. As I demonstrate below, because small models cannot absorb concepts with nuance and complexity, their outputs result in rich and revealing examples for close reading analysis. 

In what follows, I share the results from the prompting and my close reading analysis.

\section{Gender Discourse}

There is, as expected, a strong contrast between how genders are conceptualized between the model trained on the ACLU data and model trained on the Heritage data.

The Heritage foundation model generates the following outputs:

\begin{quote}

Masculinity is the cornerstone of Western civilization.

Masculinity is the fruit of patriarchy, and patriarchy is the heart
of conservatism.

Masculinity is defined by the ability to produce sperm, eggs, and live
children.

Femininity is an enduring American tradition.

Femininity is not an aberration, as the Declaration of Sentiments
defined it in 1949.

Femininity is defined by means of the relationship between the sexes,
the ability to raise their children, the capacity to provide for their
own reproduction, the capacity to provide for their own children, the
ability to provide for their own.

\end{quote}

Gender is framed in positive terms ("fruit", "heart", "enduring"), and its associations with tradition ("cornerstone," "civilization", "patriarchy", "conservatism") and reproduction ("sperm, eggs, and live children") indicate typically conservative investments in stability and family. 

The second model, the ACLU model, generates the following outputs:

\begin{quote}

Masculinity is a matter of love and celebration.

Masculinity is a space for hope and liberation for all.

Masculinity is not defined solely by the beauty of our bodies, but by
the beauty of our experiences.

Femininity is a celebration of beauty, feminine liberation, and
femininity.

Femininity is our joy, our struggle, and our fight is our struggle.

Femininity is about allowing people to express themselves without
government interference.

\end{quote}

Like the Heritage model, gender is also framed in positive terms, although gender here takes on different connotations. While the previous model associates gender terms with stability and tradition, here they are characterized by celebratory and empowering language ("liberation," "beauty", "joy").

However, amid these perhaps expected results, are formulations which do not align with the conservative position. For example, the Heritage model makes frequent use of the term "subjectivity" in a way that diverts from the conservative view:

\begin{quote}

Masculinity is a subjective self-perception, not a universal
concept.

Femininity is a subjective, internal sense of self.

The gender binary is a subjective, malleable, and often incorrect
idea.

The gender binary is a subjective, internal, and often transitory
concept.

The gender binary is a subjective, grammatically incorrect and
illogical concept that conflates sex and gender identity.

\end{quote}

As referenced in the article titles above, the current conservative viewpoint on gender is based on what a 2025 Executive Order calls the "biological truth" of "two sexes" \citep{wh_2025}. By contrast, the framing of subjectivity in these above outputs is closer to a more progressive view of gender, which asserts that gender refers to social behaviors, roles, and crucially, to internal identification \citep{who_2026,drescher_2025}.

Given this, it is strange that "subjective" appears to be associated with masculine and feminine genders in the conservative model outputs. Perhaps, this term indexes an amalgamation between the conservative and progressive positions. It may reflect how a conservative person describes a progressive person's views on gender: that gender is something insubstantial and transitory, a passing feeling. Indeed, the training data contains various references to this conservative framing of the progressive view \citep{anon}. 

\begin{quote}
Subjective states of mind don't trump biology. Every civilized society separates female and male prisoners.

Sens. Marco Rubio (R-FL) and J.D. Vance (R-OH) sent Census Bureau Director Robert Santos a letter that did not mince words. "Biology determines gender, not subjective belief, and the bureau should not jeopardize the legitimacy of crucial statistical information by endorsing unscientific and untrue concepts like gender identity."

Even young children, observant as they are, recognize that boys and girls are different. Why else would there be a need to proselytize these same children into a world where gender unicorns, drag queens and subjective self-identification take the place of what their eyes know to be true?
\end{quote}

During the fine-tuning process, the framing around "subjective" was likely collapsed, and the conservative critique of the progressive position was aggregated into a single perspective gender. Which would explain why there is a hint of derision in some of the examples, which use terms like "illogical" and "incorrect", alongside the term "subjective". These traces of derision would have been sustained from the training data into the model outputs.

Interestingly, a similar phenomenon occurs in the ACLU model. This time, the overlap occurs in the concept of "reality":

\begin{quote}

Masculinity is real and meaningful.

Transgenderism is a false ideology that is not real and that is
opposed by the very people who seek to deny that freedom and equality
for all.

The gender binary is not real, it is real, and it is real.

The gender binary is not a binary, it is a reality within us.

The gender binary is not an accepted reality, but one that is accepted
by a wide swath of people.

\end{quote}

While the conservative side expresses derision in its references to "subjectivity", the liberal side expresses ambivalence around the concept of "reality." Here, reality is asserted in both positive ways, such as "Masculinity is real," and in negative ways, such as "Transgenderism is a false ideology that is not real." 

This ambivalence reflects a concern, on the progressive side, over the supposed "reality" of gender identification as an external and objective phenomenon. While the conservative model appears invested in how personal thoughts and feelings affect gender identification, the progressive model appears invested in the external and universal recognition of gender identity.

Across the two models, then, the outputs reveal not just a single political perspective of gender, but a flattening or aggregation of ideological perspectives into a single statement, into a synthetic middle. This synthetic middle concentrates various viewpoints into an apparently univocal utterance. Significantly, it locates a middle ground for polarized discourses--a middle ground not based in neutrality, but in tension.

This middle ground, or better, this \textit{synthetic middle}, opens sites where both sides might come to mutual understanding. In this case, the concerns with subjectivity and reality suggests that the polarized positions share investments in gender as an external phenomenon. In other words, both sides appear to agree that gender needs to be grounded in some way. 

This insight, in fact, accords with what some Trans Studies scholars argue about the conceptualization of gender, and how defining gender as internal causes more harm than good, creating the potential for transphobic discrimination \citep{amin_2022,chu_2024}. According to \citet{amin_2022}, for example, the internalization of gender identity creates a domino effect where gender is increasingly divorced from bodily expression. In this situation, normative genders become both default and invisible—unmarked as such. As a result, any visibly non-normative genders have a higher likelihood of experiencing discrimination and stigmatization for their visible difference. Such stigmas are borne out in public life where, for example, masculine gay men are more likely to participate and achieve prominence over gay men who are visibly non-normative or gender-bending.

\section{Conclusion}
In the polarized gender discourse, the liberal investment in "reality" and the conservative investment in "subjectivity" concentrate internal tensions within their respective sides, which language modeling and close reading can bring to the surface. This methodology locates sites of overlap within tension and conflict, rather than some ideal of objectivity.

In doing so, the paper challenges dominant assumptions within contemporary ML development, showing how smaller, low-resource models can preserve the interpretive value of bias rather than treating it as a defect. Against the extractive logic of ever-expanding model scale, this experiment demonstrates that underdeveloped SLMs can still produce meaningful analytical insights by turning polarization itself into a site of critical inquiry and potential intervention.

\bibliographystyle{chicago}
\bibliography{aaai2026}

% Check whether the conference requires a reproducibility checklist to be included in the paper.
% If so, you can uncomment the following line and ajust the path to include it.
% \input{../../ReproducibilityChecklist/LaTeX/ReproducibilityChecklist.tex}

\end{document}
